\section{What Claude Code is especially good at}
\label{sec:strengths}

This section is a catalogue of edits that are slow, error-prone, or outright tedious by hand, and that Claude Code does well. Each example is a real prompt you can paste.

\subsection{Global notation and style changes}

\begin{lstlisting}[style=prompt]
> Replace every `\cite{...}` with `\citet{...}` where the citation is used as a grammatical noun ("X et al. showed"), and `\citep{...}` otherwise. Show me the full diff before writing.
\end{lstlisting}

\begin{lstlisting}[style=prompt]
> We have been inconsistent about the symbol for the learning rate. Pick one (prefer `\eta`), and replace `\alpha` with `\eta` everywhere it denotes learning rate (not when it denotes significance level). List any ambiguous cases for me to resolve.
\end{lstlisting}

\subsection{Cross-reference audits}

\begin{lstlisting}[style=prompt]
> Audit all cross-references in the paper. List:
>   1. every `\label` that is never `\ref`-ed,
>   2. every `\ref` whose target label does not exist,
>   3. every `\cite` key not present in references.bib.
> Do not change any files yet; just produce the report.
\end{lstlisting}

\subsection{Bibliography hygiene}

\begin{lstlisting}[style=prompt]
> In references.bib, normalise every entry: authors in "Last, First" form, journal names in title case, year present, DOI field present where known. Flag entries that look incomplete or suspicious.
\end{lstlisting}

\subsection{Length surgery}

\begin{lstlisting}[style=prompt]
> The introduction is currently ~650 words and must fit in 400. Shorten it, preserving every factual claim and every citation. Do not remove equations. Produce the diff only; I will approve.
\end{lstlisting}

\subsection{Acronym and nomenclature passes}

\begin{lstlisting}[style=prompt]
> Build an acronym list: scan the paper for capitalised acronyms, find the first occurrence of each, and ensure it is expanded on first use ("long form (LF)"). Emit a glossary table at the start of the appendix.
\end{lstlisting}

\subsection{Response-to-reviewers drafting}

\begin{lstlisting}[style=prompt]
> Reviewer 2 comment R2.4 (in review.txt) asks why we did not compare to method X. I added section 4.3 and figure 5 to address this. Draft the response entry for R2.4, citing the exact section and figure, in the tone of the existing responses in response.tex.
\end{lstlisting}

\subsection{Structural refactors}

\begin{lstlisting}[style=prompt]
> Split sections/methods.tex into three files, one per subsection, update main.tex to `\input` them, and keep the resulting PDF byte-identical as far as possible.
\end{lstlisting}

\subsection{Figure and table scaffolding}

\begin{lstlisting}[style=prompt]
> Given the numbers in results.csv, produce a booktabs table with method in the first column and the three metrics as columns 2--4, with the best value in each column in bold. Place it in sections/results.tex where the TODO marker is.
\end{lstlisting}

\subsection{British/American consistency}

\begin{lstlisting}[style=prompt]
> The paper should be British English. Find every Americanism (spellings, date formats, quotation style) and fix it. Preserve quoted material verbatim.
\end{lstlisting}

\subsection{Build verification}

\begin{lstlisting}[style=prompt]
> Run `latexmk -pdf main.tex`. If it fails, read the log, fix the cause, and re-run until it builds cleanly with no undefined references and no overfull boxes greater than 5pt.
\end{lstlisting}

\begin{tipbox}
Patterns that work well: ask for the plan or the report first, approve, then ask Claude to make the edit. You keep more control and the diffs stay small.
\end{tipbox}
